\section*{Formal Definition: Adaptive k Selection}
\label{sec:adaptive_k_formal}

\subsection*{Input Definition}

Let $D = w_1 w_2 \ldots w_n$ be a document with $n$ tokens. The document is chunked into $m$ passages:
\begin{equation}
\mathcal{C} = \{c_1, c_2, \ldots, c_m\} \quad \text{where } c_i = w_{s_i} \ldots w_{e_i} \text{ (contiguous span)}
\end{equation}

Each chunk $c_i$ is associated with metadata: document ID, position, and token span $[s_i, e_i]$.

The implicit query $q$ represents the summarization request. In practice, $q$ may be derived from the document title, first sentence, or a learned representation.

\subsection*{Retrieval Scoring Function}

After hybrid retrieval (dense + sparse + ColBERT) and cross-encoder reranking, we obtain scores for each chunk:
\begin{equation}
\text{score}_{\text{rerank}} : (q, c_i) \mapsto s_i \in \mathbb{R}
\end{equation}

These scores are sorted in descending order:
\begin{equation}
\mathbf{s} = [s_1, s_2, \ldots, s_N] \quad \text{where } s_1 \geq s_2 \geq \ldots \geq s_N
\end{equation}

Let $\sigma: \text{rank position} \to \text{chunk index}$ be the permutation, so $c_{\sigma(i)}$ is the $i$-th ranked chunk.

\subsection*{Adaptive k Selection Rule}

The adaptive k selection solves:
\begin{equation}
k^* = \arg\max_{k \in [k_{\min}, k_{\max}]} Q(k)
\end{equation}

where the quality function is:
\begin{equation}
Q(k) = \underbrace{\sum_{i=1}^{k} s_i}_{\text{cumulative confidence}} - \lambda \cdot \underbrace{\max(0, s_{k-1} - s_k)}_{\text{score dropoff}}
\end{equation}

\paragraph{Interpretation}
\begin{itemize}
\item \textbf{Cumulative confidence term}: $\sum_{i=1}^{k} s_i$ measures total relevance of top $k$ passages. Higher cumulative score indicates more content is relevant.

\item \textbf{Dropoff penalty term}: $\max(0, s_{k-1} - s_k)$ penalizes selection that includes a sharp drop in scores. This detects the boundary between high-confidence and low-confidence passages.
  \begin{itemize}
  \item If $s_{k-1} \approx s_k$ (gradual decrease), penalty $\approx 0$, allowing higher $k$.
  \item If $s_{k-1} \gg s_k$ (sharp cliff), penalty is large, encouraging lower $k$.
  \end{itemize}

\item \textbf{Hyperparameter $\lambda$}: Controls relative weight of dropoff. Higher $\lambda$ makes the algorithm more conservative (prefers fewer passages). Empirically, $\lambda = 1.0$.

\item \textbf{Constraints}: $k_{\min} \geq 1$ (minimum coverage), $k_{\max} \leq N$ (feasibility). Typically $k_{\min} = 5, k_{\max} = 30$.
\end{itemize}

\subsection*{Algorithm}

\begin{algorithm}
\caption{Adaptive k Selection}
\label{algo:adaptive_k}
\begin{algorithmic}
\Require Scores $\mathbf{s} = [s_1, \ldots, s_N]$ (sorted descending)
\Require Parameters: $k_{\min}, k_{\max}, \lambda$
\Ensure Optimal passage count $k^*$

\State $Q_{\max} \gets -\infty$
\State $k^* \gets k_{\min}$
\State $\text{cumsum} \gets 0$

\For{$k = k_{\min}$ \textbf{to} $\min(k_{\max}, N)$}
    \State $\text{cumsum} \gets \text{cumsum} + s_k$ \Comment{Running cumulative sum}
    \State $\text{dropoff} \gets \max(0, s_{k-1} - s_k)$ \Comment{Score drop at position $k$}
    \State $Q_k \gets \text{cumsum} - \lambda \cdot \text{dropoff}$
    \If{$Q_k > Q_{\max}$}
        \State $Q_{\max} \gets Q_k$
        \State $k^* \gets k$
    \EndIf
\EndFor

\State \Return $k^*$
\end{algorithmic}
\end{algorithm}

\subsection*{Complexity Analysis}

\paragraph{Time Complexity}
\begin{itemize}
\item Maintaining running cumulative sum: $O(1)$ per iteration.
\item Loop over $k \in [k_{\min}, k_{\max}]$: $O(k_{\max} - k_{\min}) = O(\Delta k)$.
\item Since $k_{\max} = 30, k_{\min} = 5$, we have $\Delta k = 25$ (constant).
\item \textbf{Total time}: $O(1)$ (independent of document length or number of retrieved passages $N$).
\end{itemize}

\paragraph{Space Complexity}
\begin{itemize}
\item Store scores $\mathbf{s}$: $O(N)$.
\item Algorithm workspace: $O(1)$.
\item \textbf{Total space}: $O(N)$.
\end{itemize}

\paragraph{Encoder and Decoder Cost}

With adaptive k, the FiD encoder cost becomes query-dependent:
\begin{equation}
\text{Cost}_{\text{encode}} = k^*(q) \cdot (\text{sequence length} \times \text{hidden dimension})
\end{equation}

Compared to fixed-k FiD:
\begin{equation}
\text{Cost}_{\text{fixed}} = k_{\text{fixed}} \cdot (\text{sequence length} \times \text{hidden dimension})
\end{equation}

\textbf{Speedup}: If $k_{\text{fixed}} = 20$ and $\mathbb{E}[k^*] = 8$, the average speedup is $\frac{20}{8} = 2.5 \times$ (60\% cost reduction).

\subsection*{Comparison with Fixed-k FiD}

\begin{table}[h]
\centering
\small
\begin{tabular}{lcc}
\toprule
\textbf{Aspect} & \textbf{Fixed-k FiD} & \textbf{Adaptive-k RALFS} \\
\midrule
Passage count $k$ & Hyperparameter (constant) & Function of $\mathbf{s}$ (variable) \\
Encoder cost per query & $O(k \cdot d)$ & $O(k^*(q) \cdot d)$ \\
Decoder cost per query & $O(k \cdot m^2)$ & $O(k^*(q) \cdot m^2)$ \\
Selection mechanism & None & Score-dropoff detection \\
Training complexity & Simple & Straight-through estimator \\
Observed speedup & 1$\times$ (baseline) & 2.5$\times$ avg. \\
Quality gain (ROUGE-2) & Baseline & +12\% \\
\bottomrule
\end{tabular}
\caption{Comparison of fixed-k FiD and adaptive-k RALFS. $d$ = hidden dimension, $m$ = sequence length.}
\label{tab:complexity}
\end{table}

\subsection*{Why Score Dropoff Works}

Retrieval rankings exhibit a characteristic pattern: top-ranked passages are relevant (high scores, tight clustering), while low-ranked passages are noise (low scores, looser distribution). The transition point between these regimes manifests as a local maximum in score difference $\delta_k = s_{k-1} - s_k$.

Our quality function $Q(k)$ implicitly detects this cliff:
\begin{itemize}
\item When scores decrease uniformly (no cliff), the dropoff penalty is small, allowing higher $k$.
\item When a sharp cliff appears, the penalty becomes large, discouraging inclusion of the cliff position.
\item The algorithm naturally selects the point just before the cliff.
\end{itemize}

This observation is supported by cascade retrieval literature (Bruch et al., 2021), which shows that stopping when ranking quality degrades yields both efficiency and quality gains.

\subsection*{Training: Straight-Through Estimator}

During training, adaptive k selection is deterministic but non-differentiable (due to $\arg\max$). We use a \emph{straight-through estimator}:

\begin{itemize}
\item \textbf{Forward pass}: Compute $k^*(q)$ deterministically using Algorithm~\ref{algo:adaptive_k}.
\item \textbf{Backward pass}: Treat $k^*$ as a constant (fixed per query). Gradients flow through the decoder to the encoder representations, but not through the k-selection mechanism itself.
\item \textbf{Effect}: The model learns to adapt its internal representations to different passage counts across different queries.
\end{itemize}

This avoids the complexity of learned k selection (which would require additional parameters and loss terms) while allowing the model to exploit variable-length passage sets.

\subsection*{Hyperparameter Sensitivity}

\paragraph{$\lambda$ (dropoff penalty weight)}
\begin{itemize}
\item $\lambda = 0$: Adaptive k reduces to cumulative score maximization; ignores dropoff.
\item $\lambda = 1.0$: Balanced weighting (empirically optimal).
\item $\lambda > 1$: Conservative; prefers fewer passages (more sensitive to drops).
\end{itemize}

\paragraph{$k_{\min}, k_{\max}$ (bounds)}
\begin{itemize}
\item $k_{\min} = 5$: Ensures minimum context for coherent generation.
\item $k_{\max} = 30$: Bounds worst-case computation.
\item Typical range: $k_{\min} \in [3, 10], k_{\max} \in [20, 50]$ depending on inference budget.
\end{itemize}

\subsection*{Extensions and Future Work}

\paragraph{Learned k Selection}
Instead of score-dropoff, a learned selector could predict $k$ from query features:
\begin{equation}
k^* = \text{MLP}_\theta(q, \mathbf{s}) \quad \text{(learned model)}
\end{equation}

Benefits: Potentially better modeling of query-$k$ dependencies. Drawback: Requires additional training data/labels.

\paragraph{Adaptive Fusion Weights}
Currently, RRF uses fixed weights. Query-adaptive fusion could select different weight combinations per query:
\begin{equation}
\alpha(q), \beta(q), \gamma(q) = \text{MLP}_\phi(q) \quad \text{(learned per query)}
\end{equation}

\paragraph{Joint Optimization}
End-to-end training of k selection and generation jointly could improve both:
\begin{equation}
\min_{\theta} \mathbb{E}_{(d,s_{\text{ref}})} [\mathcal{L}_{\text{gen}}(s_{\text{ref}}, \text{GEN}(\mathcal{R}_{k^*}, q)) - \alpha \cdot k^*]
\end{equation}

This trades off generation quality and computational cost explicitly, potentially converging to a better equilibrium.
